\subsection{Generating Modules and Finitely Generated Projective Modules}

PROPOSITION 2. --- Let A and B be algebras over k, and let P be an \((A,B)_{k}\) -bimodule. Suppose that the mapping \(b \mapsto b_{P}\) from B to \(\operatorname{End}_{\mathrm{A}}(\mathrm{P})\) is bijective. The following assertions are equivalent:

\begin{enumerate}
\def\labelenumi{(\roman{enumi})}
\item
  The A-module P is projective and finitely generated.
\item
  The mapping \(\Theta\) (VIII, p.~96) is an isomorphism of \((\mathrm{B},\mathrm{B})_k\) -bimodules from \(\mathrm{P}^* \otimes_{\mathrm{A}} \mathrm{P}\) to \(s\mathrm{B}_d\) .
\item
  The image of \(\Theta\) contains the unit element of B.
\end{enumerate}

If, moreover, the mapping \(a \mapsto a_{P}\) from A to \(\operatorname{End}_{\mathrm{B}}(\mathrm{P})\) is bijective, then the assertions above are equivalent to the following statement:

\begin{enumerate}
\def\labelenumi{(\roman{enumi})}
\setcounter{enumi}{3}
\tightlist
\item
  There exist a \((\mathrm{B},\mathrm{A})_{k}\) -bimodule Q and a surjective homomorphism of \((\mathrm{B},\mathrm{B})_{k}\) -bimodules from \(Q\otimes_{A}P\) to \(_{s}B_{d}\) .
\end{enumerate}

By the corollary of II, §4, No.~2, p.~271, assertion (i) implies (ii). Moreover, (iii) follows from (ii). Let \(t \in \mathrm{B}^* \otimes_{\mathrm{A}} \mathrm{B}\) be such that \(\Theta(t) = 1\) . Let \(n\) be an integer, and let \((x_1, \ldots, x_n) \in \mathrm{P}^n\) and \((x_1^*, \ldots, x_n^*) \in \mathrm{P}^{*n}\) be such that \(t = \sum_{i=1}^{n} x_i^* \otimes x_i\) . For every \(x\) belonging to \(\mathrm{P}\) , the relation \(x\Theta(t) = x\) can be written as

\[
\sum_ {i = 1} ^ {n} \langle x, x _ {i} ^ {*} \rangle x _ {i} = x.
\]

It follows that the A-module P is finitely generated by the family \((x_{i})_{1\leqslant i\leqslant n}\) , and we conclude the proof of the implication (iii) \(\Rightarrow\) (i) using Proposition 12 of II, §2, No.~6, p.~238.

We obviously have (ii) \(\Rightarrow\) (iv). Let \(Q\) be a \((B,A)_k\) -bimodule and \(\theta\) a surjective homomorphism of \((B,B)_k\) -bimodules from \(Q \otimes_A P\) to \(s B_d\) . In the notation of the previous subsection, there exists a homomorphism of \((B,A)_k\) -bimodules \(\zeta: Q \to \widetilde{P}\) such that \(\theta(y \otimes x) = \langle \zeta(y), x \rangle\) for \(x \in P\) and \(y \in Q\) , and we have \(\theta = \widetilde{\Lambda} \circ (\zeta \otimes 1_P)\) . Since \(\theta\) is surjective, the same holds for \(\widetilde{\Lambda}\) . If the mapping \(a \mapsto a_P\) from \(A\) to \(\mathrm{End}_B(P)\) is bijective, then \(\Theta\) is surjective by relation (14) of Proposition 1 of VIII, p.~97; the implication (iv) \(\Rightarrow\) (iii) follows.

PROPOSITION 3. --- Let A and B be algebras over k and P an \((A,B)_{k}\) -bimodule. The following properties are equivalent:

\begin{enumerate}
\def\labelenumi{(\roman{enumi})}
\item
  The A-module P is generating.
\item
  The image of the mapping \(\Lambda\) from \(\mathrm{P} \otimes_{\mathrm{B}} \mathrm{P}^{*}\) to \(s\mathrm{A}_{d}\) (VIII, p.~95) contains the unit element.
\item
  There exist a \((\mathrm{B},\mathrm{A})_{k}\) -bimodule Q and a surjective homomorphism of \((\mathrm{A},\mathrm{A})_{k}\) -bimodules from \(P\otimes_{B}Q\) to \(_{s}A_{d}\) .
\end{enumerate}

If, moreover, the mapping \(b \mapsto b_{P}\) from B to \(\operatorname{End}_{\mathrm{A}}(\mathrm{P})\) is bijective, then the assertions are equivalent to the following statement:

\begin{enumerate}
\def\labelenumi{(\roman{enumi})}
\setcounter{enumi}{3}
\tightlist
\item
  The mapping \(\Lambda\) is an isomorphism from \(\mathrm{P} \otimes_{\mathrm{B}} \mathrm{P}^{*}\) to \(s\mathrm{A}_d\) .
\end{enumerate}

(i) \(\Leftrightarrow\) (ii): The image of \(\Lambda\) is the trace ideal \(\tau(\mathrm{P})\) (VIII, p.~80). The equivalence of (i) and (ii) therefore follows from Theorem 1 of VIII, p.~80.

(ii) \(\Rightarrow\) (iii): It suffices to take \(Q = P^{*}\) .

(iii) \(\Rightarrow\) (ii): Let \(Q\) be a \((B,A)_k\) -bimodule and \(\psi\) a homomorphism of \((A,A)_k\) -bimodules from \(P \otimes_B Q\) to \({}_s A_d\) . There exists a homomorphism of \((B,A)_k\) -bimodules \(\Psi\) from \(Q\) to \(P^*\) such that \(\psi(x \otimes y) = \langle x, \Psi(y) \rangle\) for \(x \in P\) and \(y \in Q\) , and we have the equality \(\psi = \Lambda \circ (1_P \otimes \Psi)\) . If \(\psi\) is surjective, then so is \(\Lambda\) .

It is clear that (iv) implies (ii). Conversely, suppose that property (ii) holds and that the mapping \(b \mapsto b_{P}\) from B to \(\operatorname{End}_{\mathrm{A}}(\mathrm{P})\) is bijective. Let e be an element of \(P \otimes_{B} P^{*}\) such that \(\Lambda(e) = 1\) . By relation (7) of VIII, p.~96, we have \(u = \Lambda(u)e\) for every u in \(P \otimes_{B} P^{*}\) ; the injectivity of \(\Lambda\) follows.

PROPOSITION 4. --- Let A and B be k-algebras and P an \((\mathrm{A},\mathrm{B})_{k}\) -bimodule. Suppose that the mapping \(b\mapsto b_{P}\) from B to \(\operatorname{End}_{\mathrm{A}}(\mathrm{P})\) is bijective.

\begin{enumerate}
\def\labelenumi{\alph{enumi})}
\item
  If the A-module P is generating, then the right B-module P is projective and finitely generated.
\item
  If the A-module P is projective and finitely generated, then the right B-module P is generating.
\end{enumerate}

Suppose that the A-module P is generating. It is then faithful and balanced (VIII, p.~82, Theorem 2), and consequently the mapping \(a \mapsto a_{\mathrm{P}}\) from A to \(\mathrm{End}_{\mathrm{B}}(\mathrm{P})\) is bijective (VIII, p.~97, Remark 1). Moreover, the mapping \(\Lambda : \mathrm{P} \otimes_{\mathrm{B}} \mathrm{P}^{*} \to {}_{s} \mathrm{A}_{d}\) is bijective (VIII, p.~98, Proposition 3). We view P as a \((\mathrm{B}^{\circ}, \mathrm{A}^{\circ})_{k}\) -bimodule. The mapping \(\Lambda\) induces a bijective mapping \(\mathrm{P}^{*} \otimes_{\mathrm{B}^{\circ}} \mathrm{P} \to \mathrm{A}^{\circ}\) ; by Proposition 2 of VIII, p.~98, (iv) \(\Rightarrow\) (i), the right B-module P is projective and finitely generated.

Now, suppose that the A-module P is projective and finitely generated. Then the mapping \(\Theta: P^{*} \otimes_{A} P \to_{s} B_{d}\) is bijective (loc. cit., (i) \(\Rightarrow\) (ii)). By the implication (iii) \(\Rightarrow\) (i) of Proposition 3 above applied to the \((\mathrm{B}^{\circ}, \mathrm{A}^{\circ})_{k}\) -bimodule P, the right B-module P is generating.

COROLLARY 1. --- The countermodule of a generating module is projective and finitely generated. The countermodule of a finitely generated projective module is generating.

Let A be a k-algebra, and let M be an A-module. We denote the opposite k-algebra of the algebra \(\operatorname{End}_{\mathrm{A}}(\mathrm{M})\) by B. The corollary follows from Proposition 4 applied to the \((\mathrm{A},\mathrm{B})_{k}\) -bimodule M.

COROLLARY 2. --- Let A and B be k-algebras and P an \((A,B)_{k}\) -bimodule. The following properties are equivalent:

\begin{enumerate}
\def\labelenumi{(\roman{enumi})}
\item
  The A-module P is generating, and the mapping \(b \mapsto b_{P}\) from B to \(\operatorname{End}_{\mathrm{A}}(\mathrm{P})\) is bijective.
\item
  The right B-module P is projective, finitely generated, faithful, and balanced, and the mapping \(a \mapsto a_{\mathrm{P}}\) from A to \(\operatorname{End}_{\mathrm{B}}(\mathrm{P})\) is bijective.
\end{enumerate}

The implication \((i) \Rightarrow (ii)\) follows from Proposition 4, a) and Remark 1 (VIII, p.~97). Assume that (ii) holds; then the A-module P is generating (Proposition 4, b) applied to the \((B^{\circ}, A^{\circ})_{k}\) -bimodule P). Since the B-module P is faithful and balanced, the second assertion of (i) is also satisfied (VIII, p.~97, Remark 1).

